\subsection{Elastic Weight Consolidation (EWC)}

\gls{ewc} \cite{ll_ewc_2017} is a regularisation-based \gls{cl} method that mitigates \gls{cf} by penalising changes to parameters that were important for previous tasks. In contrast to replay-based methods, \gls{ewc} does not store raw past samples; instead, it carries forward a compact summary of past tasks through parameter snapshots at task transitions and importance weights based on curvature of the loss landscape.

\paragraph{Parameter Importance via Fisher Information}
The Fisher information matrix provides a principled measure of parameter importance: parameters whose perturbation strongly affects the model's output distribution are assigned higher importance and are more strongly penalised by \gls{ewc}. The Fisher information matrix $F \in \mathbb{R}^{|\theta| \times |\theta|}$ characterises the curvature of the log-likelihood surface with respect to the model parameters. Its entries are defined as:
%
\begin{equation}\label{eqn:ewc_fisher_full}
    F_{ij} = \mathbb{E}_{(x,y)\sim\mathcal{D}}\left[
        \frac{\partial \log p_\theta(y \mid x)}{\partial \theta_i} \cdot
        \frac{\partial \log p_\theta(y \mid x)}{\partial \theta_j}
    \right]
\end{equation}
%
Each entry $F_{ij}$ captures the interaction between parameters $\theta_i$ and $\theta_j$ with respect to the model's output distribution: large off-diagonal entries indicate that the two parameters are statistically coupled, such that moving either one affects the informativeness of the other. For a model with $|\theta|$ parameters, storing $F$ exactly requires $\mathcal{O}(|\theta|^2)$.

Rather than computing the full $|\theta| \times |\theta|$ matrix, \gls{ewc} approximates it by its diagonal (i.e. discarding parameter coupling), reducing storage to $\mathcal{O}(|\theta|)$, where $i$ indexes each scalar element of the full parameter vector $\theta \in \mathbb{R}^{|\theta|}$. The diagonal entry for parameter $\theta_i$ after task $t$ is:
%
\begin{equation}\label{eqn:ewc_fisher}
    F_{t,i} \approx \mathbb{E}_{(x,y)\sim\mathcal{D}_t}
    \left[
        \left(
            \frac{\partial \log p_\theta(y\mid x)}{\partial \theta_i}
        \right)^2
    \right]
\end{equation}
%
where $\mathcal{D}_t$ is the data distribution for task $t$, $p_\theta(y \mid x)$ is the model's predicted class probability, and $(x, y)$ denotes an input--label pair drawn from $\mathcal{D}_t$.

In practice, this expectation is approximated using squared gradients of the training loss accumulated over minibatches. For a minibatch of size $B$, running sums $S_{t,i}$ and sample count $N_t$ are maintained as:
%
\begin{equation}\label{eqn:ewc_fisher_minibatch}
    S_{t,i} \leftarrow S_{t,i} + B\, g_i^2,
    \qquad
    N_t \leftarrow N_t + B
\end{equation}
%
where $g_i = \frac{\partial \mathcal{L}}{\partial \theta_i}$ is the gradient of the classification loss $\mathcal{L}$ with respect to $\theta_i$. At the end of task $t$, the per-task Fisher estimate is the sample-weighted average:
%
\begin{equation}\label{eqn:ewc_fisher_taskend}
    \hat{F}_{t,i} = \frac{S_{t,i}}{N_t}
\end{equation}

Rather than storing a separate Fisher tensor per task, which would incur memory cost linear in the number of tasks, a single diagonal Fisher tensor $F_i$ is maintained and updated incrementally. After the first task, $F_i \leftarrow \hat{F}_{t,i}$. Upon completing each subsequent task, the new estimate is merged into the running importance via arithmetic averaging:
%
\begin{equation}\label{eqn:ewc_fisher_merge}
    F_i \leftarrow \frac{k\, F_i + \hat{F}_{t,i}}{k+1},
    \qquad k \leftarrow k+1
\end{equation}
%
where $k$ is the number of tasks already incorporated into $F_i$ prior to the update. The quadratic penalty therefore uses this merged $F_i$ alongside a single set of reference parameters $\theta^* \in \mathbb{R}^{|\theta|}$, overwritten at each task boundary with the current weights. This differs from the original formulation \cite{ll_ewc_2017}, which retains every task optimum $\theta^*_t$ and sums $\sum_t F_{t,i}(\theta_i - \theta_{t,i}^*)^2$, but achieves the same regularisation effect at constant memory cost regardless of the number of tasks.

\paragraph{\gls{cl} Approach}

When learning a new task, \gls{ewc} augments the standard classification loss with a quadratic penalty that pulls parameters back toward previously consolidated values:
%
\begin{equation}\label{eqn:ewc_loss}
    \mathcal{L}_\text{EWC}
    =
    \mathcal{L}_\text{CE}
    +
    \frac{\lambda}{2}
    \sum_i F_i\,(\theta_i-\theta_i^*)^2
\end{equation}
%
where $\mathcal{L}_\text{CE}$ is the current-task cross-entropy, $\lambda$ controls the stability--plasticity trade-off, and $(F_i,\theta_i^*)$ denote the consolidated importance and reference parameters from prior tasks. This objective allows parameters with low estimated importance to adapt more freely, while constraining parameters that are critical to old tasks, as illustrated in Fig.~\ref{fig:ll_ewc}.

This is effective because the Fisher acts as a local curvature proxy of the old-task objective around $\theta^*$. Parameters with high curvature correspond to directions where small perturbations strongly change old-task predictions, so the quadratic penalty enforces stability in those directions. Conversely, flat directions (small Fisher entries) remain plastic, enabling adaptation to new tasks. 
% Under a Laplace approximation, this can be interpreted as sequential Bayesian posterior updating with a Gaussian prior centred at $\theta^*$ and precision given by the Fisher.

% \begin{equation}
%     p(\theta~|~\mathcal{D}_t) \approx \mathcal{N}(\theta^*, F^{-1})
% \end{equation}

\paragraph{Inference}
\gls{ewc} uses a standard linear classifier at inference time. In \gls{til}, task identity is used to select the relevant task logit slice. In \gls{cil}, prediction is made from the full logit vector across all seen classes.